% =====================================================================
%  Multi-annotator fusion methods
%  Drop into a methods chapter. Requires: amsmath, booktabs, algorithm,
%  algpseudocode (or algorithm2e; pseudo-code is plain enough to swap).
% =====================================================================

\section{Multi-annotator fusion}
\label{sec:annotator-fusion}

The real validation stack (033) ships with three independently produced
point-annotation files; each annotator marks the centre of every CCP they see.
To build a single training mask the per-frame point sets
$\mathcal{P}_1, \mathcal{P}_2, \mathcal{P}_3 \subset \mathbb{R}^2$
must be fused into one set of centres
$\mathcal{C} \subset \mathbb{R}^2$, with the property that two points labelled
by different annotators within a small radius~$\tau$ collapse to the same
centre (they are interpreted as the same physical CCP), while points farther
apart remain separate. The chosen $\tau$ acts as the operational definition
of ``same CCP''. Two clustering procedures were evaluated for this step.

% ---------------------------------------------------------------------
\subsection{Greedy single-pass clustering}
\label{sec:greedy-fusion}

The greedy procedure walks through the pooled annotator points in some order
and, for each point~$p$, either joins the nearest existing cluster (if its
running centroid lies within~$\tau$) or starts a new singleton cluster.
The centroid is updated on the fly as the running mean of all members.
The full procedure is given as Algorithm~\ref{alg:greedy}.

\begin{algorithm}[ht]
\caption{Greedy fusion of per-annotator points}
\label{alg:greedy}
\begin{algorithmic}[1]
\Require ordered sequence of points $p_1, \dots, p_N$; threshold $\tau$
\State $\mathcal{K} \gets \emptyset$ \Comment{set of clusters, each a list of points}
\For{$p$ in $p_1, \dots, p_N$}
  \If{$\mathcal{K} = \emptyset$}
    \State add new cluster $\{p\}$ to $\mathcal{K}$; \textbf{continue}
  \EndIf
  \State $c^\star \gets \arg\min_{c \in \mathcal{K}} \lVert \bar{c} - p\rVert_2$
         \Comment{$\bar{c}$ is the running centroid}
  \If{$\lVert \overline{c^\star} - p\rVert_2 < \tau$}
    \State append $p$ to cluster $c^\star$;\ \ update $\overline{c^\star}$
  \Else
    \State add new cluster $\{p\}$ to $\mathcal{K}$
  \EndIf
\EndFor
\State \Return $\{\, \overline{c} : c \in \mathcal{K} \,\}$
\end{algorithmic}
\end{algorithm}

\paragraph{Properties.}
Complexity is $\mathcal{O}(N K)$ per frame, where $K$ is the number of
clusters discovered. The output depends on the visiting order: with three
annotators, deterministically processing annotator~1 first would bias every
cluster toward that annotator's coordinate convention. To remove the bias
while keeping runs reproducible, the per-frame annotator order is shuffled
using a deterministic random permutation seeded by the frame index
(``random per-frame shuffle''). Crucially, the running-centroid update with
a hard $\tau$ on the \emph{nearest-centroid} test admits chain merging:
three points that are pairwise within $\tau$ of \emph{some} chain neighbour
may collapse into a single cluster whose maximum internal distance exceeds
$\tau$. Empirically this acts as a denoiser when three annotators
disagree by a few pixels on the location of the same physical CCP.

% ---------------------------------------------------------------------
\subsection{Agglomerative complete-linkage clustering}
\label{sec:complete-fusion}

The complete-linkage method (\texttt{scipy.cluster.hierarchy.linkage} with
\texttt{method="complete"}, followed by \texttt{fcluster} at threshold $\tau$)
starts from singleton clusters and repeatedly merges the pair $(A, B)$
whose \emph{maximum} pairwise distance is smallest:
\[
  d_{\text{complete}}(A, B) \;=\;
  \max_{a \in A,\ b \in B} \lVert a - b \rVert_2,
\]
stopping as soon as no further merge satisfies $d_{\text{complete}}(A,B) < \tau$.
Equivalently, the procedure returns the coarsest partition such that the
\emph{diameter} (maximum within-cluster pairwise distance) of every cluster
is strictly below~$\tau$.

\paragraph{Properties.}
Complexity is $\mathcal{O}(N^2 \log N)$ per frame and $\mathcal{O}(N^2)$
memory for the pairwise distance matrix; tractable for the
$\mathcal{O}(100)$ points per frame in the present data. The output is
\emph{order-independent} by construction. Unlike the greedy procedure,
the threshold $\tau$ is a hard cap on cluster diameter (no chain merging).

% ---------------------------------------------------------------------
\subsection{Empirical comparison}
\label{sec:fusion-comparison}

Both methods were evaluated by training 5-seed grids on the real CCP data
under the contiguous train/val split (\S\ref{??}). Table~\ref{tab:fusion-deta}
reports the early-stopping val DetA averaged over annotators 1--3 for the
full $2 \times 3$ grid of (algorithm, threshold) combinations. The
corresponding per-seed values are listed in
Tables~\ref{tab:fusion-deta-val-perseed} and~\ref{tab:fusion-deta-test-perseed}.

\begin{table}[ht]
\centering
\caption{Validation DetA under different multi-annotator fusion choices,
mean $\pm$ standard deviation across 5 random seeds. Rows isolate the
algorithm choice ($\tau = 5$ fixed); columns isolate the threshold choice
(complete-linkage fixed).}
\label{tab:fusion-deta}
\begin{tabular}{l l c c c}
\toprule
 &  & \textbf{Greedy} & \multicolumn{2}{c}{\textbf{Complete linkage}} \\
\cmidrule(lr){3-3}\cmidrule(lr){4-5}
Model    & & $\tau{=}5$           & $\tau{=}5$           & $\tau{=}4$           \\
\midrule
UNet     & & $0.7783 \pm 0.0118$  & $0.7620 \pm 0.0127$  & $0.7463 \pm 0.0069$  \\
EquiUNet & & $0.7831 \pm 0.0059$  & $0.7855 \pm 0.0040$  & $0.7685 \pm 0.0065$  \\
\bottomrule
\end{tabular}
\end{table}

\paragraph{Threshold effect (architecture-independent).}
Lowering $\tau$ from $5$ to $4$ under complete-linkage costs UNet
$\Delta \mathrm{DetA} = -0.016$ and EquiUNet $\Delta \mathrm{DetA} = -0.017$
--- essentially identical degradations. The tighter diameter cap splits
some physical CCPs (annotator disagreement up to $\sim 6$\,px around the
true centre) into two close centres that the network must then detect
separately, producing a denser and noisier training target irrespective of
architecture.

\paragraph{Algorithm effect (architecture-dependent).}
At a fixed $\tau = 5$, complete-linkage versus greedy gives
$\Delta \mathrm{DetA} = -0.016$ on UNet but only $+0.002$ on EquiUNet
--- statistically indistinguishable from zero given the per-seed standard
error~$\approx 0.005$. UNet pays a measurable cost for the extra centres
that complete-linkage refuses to merge; EquiUNet does not.

\paragraph{Interpretation.}
Greedy fusion's running-centroid update admits chain-merging: three
annotator points within pairwise reach of \emph{some} chain neighbour
collapse into a single centre even when their pairwise diameter exceeds
$\tau$. Empirically, this acts as a denoiser when three annotators disagree
by a few pixels on the same physical CCP. Complete-linkage, by enforcing a
hard diameter cap, instead splits such cases into two close centres.
For UNet the extra centres hurt because the network must learn to detect
both as separate peaks; for EquiUNet, the built-in rotational symmetry
constraint reduces the effective hypothesis space enough that the noisier
target does not translate into a measurable loss --- the network's
inductive bias absorbs the upstream noise. The greedy procedure with
$\tau = 5$ was retained as the operational fusion method because it is
unambiguously best for UNet and on par with the best EquiUNet
configuration, providing a single default that is robust across
architectures.

\paragraph{Per-seed breakdown.}
The aggregate numbers in Table~\ref{tab:fusion-deta} are backed by the
per-seed values in Table~\ref{tab:fusion-deta-val-perseed} (validation) and
Table~\ref{tab:fusion-deta-test-perseed} (held-out 012 test stack). Test
results follow the same qualitative pattern as validation, with one
caveat: the test stack is a single 46-frame recording, so per-seed noise
is larger ($\sigma_\text{test} \approx 2 \sigma_\text{val}$), and the
absolute test--val difference reflects domain shift between the validation
and test cells rather than any artefact of the fusion choice.

\begin{table}[ht]
\centering
\caption{Validation DetA per seed for the six fusion configurations. Each
row reports the early-stopping val DetA (mean over the three annotator
CSVs) for the five training seeds, plus the across-seed mean $\pm$
standard deviation.}
\label{tab:fusion-deta-val-perseed}
\begin{tabular}{l l c ccccc c}
\toprule
Model & Algo & $\tau$ & seed 0 & seed 1 & seed 2 & seed 3 & seed 4 & Mean $\pm$ Std \\
\midrule
UNet     & Greedy   & 5 & 0.7778 & 0.7727 & 0.7699 & 0.7987 & 0.7724 & $0.7783 \pm 0.0118$ \\
UNet     & Complete & 5 & 0.7624 & 0.7789 & 0.7473 & 0.7692 & 0.7522 & $0.7620 \pm 0.0127$ \\
UNet     & Complete & 4 & 0.7534 & 0.7515 & 0.7412 & 0.7483 & 0.7372 & $0.7463 \pm 0.0069$ \\
EquiUNet & Greedy   & 5 & 0.7841 & 0.7871 & 0.7868 & 0.7727 & 0.7846 & $0.7831 \pm 0.0059$ \\
EquiUNet & Complete & 5 & 0.7869 & 0.7844 & 0.7837 & 0.7914 & 0.7809 & $0.7855 \pm 0.0040$ \\
EquiUNet & Complete & 4 & 0.7794 & 0.7687 & 0.7627 & 0.7671 & 0.7649 & $0.7685 \pm 0.0065$ \\
\bottomrule
\end{tabular}
\end{table}

\begin{table}[ht]
\centering
\caption{Test DetA per seed on the held-out 012 stack. Same six fusion
configurations as Table~\ref{tab:fusion-deta-val-perseed}.}
\label{tab:fusion-deta-test-perseed}
\begin{tabular}{l l c ccccc c}
\toprule
Model & Algo & $\tau$ & seed 0 & seed 1 & seed 2 & seed 3 & seed 4 & Mean $\pm$ Std \\
\midrule
UNet     & Greedy   & 5 & 0.8147 & 0.8142 & 0.7683 & 0.8172 & 0.8210 & $0.8071 \pm 0.0218$ \\
UNet     & Complete & 5 & 0.7812 & 0.8143 & 0.8165 & 0.8196 & 0.7850 & $0.8033 \pm 0.0186$ \\
UNet     & Complete & 4 & 0.7702 & 0.8115 & 0.7680 & 0.7874 & 0.6902 & $0.7655 \pm 0.0455$ \\
EquiUNet & Greedy   & 5 & 0.8229 & 0.8027 & 0.8042 & 0.7964 & 0.8317 & $0.8116 \pm 0.0150$ \\
EquiUNet & Complete & 5 & 0.8131 & 0.7829 & 0.7814 & 0.8057 & 0.8146 & $0.7995 \pm 0.0162$ \\
EquiUNet & Complete & 4 & 0.8304 & 0.7806 & 0.7855 & 0.7627 & 0.7936 & $0.7906 \pm 0.0250$ \\
\bottomrule
\end{tabular}
\end{table}
